\documentclass[11pt]{article}
\usepackage[margin=1in]{geometry}

\begin{document}
\pagestyle{empty}

\begin{center}
\textbf{Significance Statement}\\[0.3em]
\textit{Displacement-Engineered Warp Fields: Tic-Tac Geometry, Stability Surfaces, and Identity-Space Confinement}
\end{center}

\vspace{0.5em}

Publicly reported aerospace observations of smooth, capsule-like craft exhibiting extreme kinematic profiles have no accepted mathematical framework for evaluation, leaving such reports outside the reach of rigorous critique. This work supplies one: a displacement-based geometric model, built on identity-space formalism, that represents the reported ``tic-tac'' geometry as a well-defined engineering object subject to explicit curvature, torsion, and stability constraints. Rather than proposing a new propulsion mechanism or stress-energy source, the model converts a qualitative observational puzzle into a falsifiable mathematical specification, complete with governing equations, a flight-envelope inequality, and engineering parameter tables. This provides researchers in geometric field theory, mathematical physics, and aerospace engineering a concrete structure to test, critique, and extend, and offers a template for applying geometric field methods to unconventional motion signatures more broadly.

\end{document}
